\scanpage{171}
\mspage{85}
% L01
ponding to this initial intuition...

% L02--L17: passage cancelled by long diagonal strokes; local cancellations retained.
\deleted{Critique of the notations. Whereas the \deleted{\unsure{letters}}
\unsure{graphic forms} $B_U$, $B_X$ are constructed in
terms of the letters $U,X$ (\unsure{suggesting|being concerned with}, \unsure{in}
\unsure{this instance|the origin}, the \unsure{schemes} $U,X$) in the
same way, they have
quite different meanings --- $B_U$ denotes the
\unsure{classifying} topos $B_{\unsure{\Pi_1|\pi_1},U}$, but $B_X$ a
distinctly bigger topos. So I prefer
to write \deleted{\illegible{1}} \unsure{finally} $B_{X;S}$ (to suggest
the sheaves on $X$ \unsure{set|essentially} constant on
$X\setminus S$ and on $S$) \deleted{\unsure{or better}, $B_{U,S}$ (to
\unsure{suggest} a topos obtained by gluing
$U$ and $S$, likewise $B$} \inserted{and likewise $B_{D;S}$} ; on the other hand
in $B_{\unsure{\Pi_1|\pi_1},D}$ and $B_{\unsure{\Pi_1|\pi_1},X}$ one would write
simply (as for $B_U$)
\inserted{$B_D$, $B_X$}}

% L18--L23
All these operations are essentially
trivial and without mystery, and independent of
\unsure{all|the} special hypotheses of the ``groups with
loops'' type. A particular interest attaches
to the case where \unsure{$B_S$|$B_D$} is a discrete topos i.e.
identifies with $B_I$, \deleted{\unsure{where}} $I=\pi_0(B_{D^*})$, so

% P: handwritten numeral in the lower-left corner; hand not identified.
\margin{170}

\scanpage{172}
% L01--L07
that we start simply with a morphism
of multi-Galois topoi
\[
 B_{D^*}\longrightarrow B_U
\]
but where in addition we have a group $\Gamma$
\inserted{(discrete, say --- or profinite in the \unsure{arithmetic}
\unsure{context})}
which operates on $B_{D^*}$, $B_U$ ; the preceding
morphism commuting with the action of $\Gamma$.

% L08--L13
Let us note that the datum of an action
of $\Gamma$ on a topos $B$ makes it possible to
construct a topos $(B,\Gamma)$, and
a morphism $(B,\Gamma)\longrightarrow B_\Gamma$
\inserted{(classifying topos of $\Gamma$}
\uuline{i.e.} a $\Gamma$-torsor of $(B,\Gamma)$ --- and

% L14--L23
$B$ identifies with the topos induced by
$\mathcal X=(B/\Gamma)$ on this $\Gamma$ torsor. Conversely,
the datum of a topos $\mathcal X$ and of a
$\Gamma$-torsor $U$ of $\mathcal X$, and of an
\unsure{isomorphism} of $B$ with the induced topos
$\mathcal X_{/U}$, \deleted{\unsure{by}} (identified with $U$, or with $B$) allows
one to recover operations of $\Gamma$ on $B$,
via the operations on $U$. Thus, the \uline{datum}

% P: handwritten numeral in the lower-left corner; hand not identified.
\margin{171}

\scanpage{173}
\mspage{86}
% L01--L08
 of an operation of $\Gamma$ on a topos \inserted{$B$} amounts
to that of the datum of $B$ as a
Galois covering with group $\Gamma$ of
another topos (\unsure{essential\textsuperscript{ly}} unique, then denoted
$(B,\Gamma)$ ... ). Thus, to have
\deleted{\unsure{$B_{D^*}$}} $\Gamma$ operate on $B_{D^*}\longrightarrow B_U$ is
the same thing as inserting this
\deleted{\unsure{diagram}} \inserted{arrow} into a diagram like

% D01
\diagram{\begin{tikzcd}[ampersand replacement=\&]
 B_{D^*}\text{\deleted{\unsure{,\Gamma}}} \arrow[r] \arrow[d] \& B_U \arrow[d] \\
 B_{D^*,\Gamma} \arrow[r] \& B_{U,\Gamma}
\end{tikzcd}}

% L09--L14
 where the vertical arrows are $\Gamma$-Galois
and the square is $2$-Cartesian, or again
(independently of the prior datum
of $B_{D^*}$, $B_U$) it is giving oneself an
\unsure{essential\textsuperscript{ly}} commutative triangle of morphisms
of topoi

% D02
\diagram{\begin{tikzcd}[ampersand replacement=\&,column sep=2.8em,row sep=2.6em]
B_{D^*,\Gamma} \arrow[rr] \arrow[dr] \& {} \& B_{U,\Gamma} \arrow[dl,no head] \\
{} \& B_{\Gamma} \& {}
\end{tikzcd}}

% P: handwritten numeral in the lower-left corner; hand not identified.
\margin{172}

\scanpage{174}
% L01--L07
If $\Gamma$ operates on a multi-Galois topos,
we see that $(B,\Gamma)$ is also
multi-Galois, and the situation of a
multi-Galois topos $B_U$ and of an
operation of $\Gamma$ on it amounts to that
of a multi-Galois topos \inserted{$B_{U,\Gamma}$} and of
a morphism $B_{U,\Gamma}\unsure{\longrightarrow|\longleftrightarrow}B_\Gamma$.

% L08--L12
So the datum of a situation
$B_{D^*}\longrightarrow B_U$ of multi-Galois topoi,
and of an operation of $\Gamma$ on it,
amounts \unsure{exactly} to that of \unsure{homomorphisms}
of multi-Galois topoi
\[
 B_{D^*,\Gamma}\longrightarrow B_{U,\Gamma}\longrightarrow B_\Gamma
\]

% L13--L16
\unsure{\deleted{The connectedness}|\uline{The connectedness}}
$B_{U,\Gamma}$ is \unsure{$0$-connected} iff $\operatorname{card}\pi_0(\unsure{B_U})/\Gamma=1$,
i.e. $B_U$ nonempty and $\Gamma$ transitive on
$\pi_0(B_U)$. --- Let us note that the situation

% P: handwritten numeral in the lower-left corner; hand not identified.
\margin{173}

\scanpage{175}
\mspage{87}
% L01--L11
 envisaged at the beginning, with $(X,S)$ over a field $K$,
whence \deleted{\unsure{$B$}} $B_{D^*}\longrightarrow B_U\longrightarrow B_\Gamma$
\inserted{``$e_{\text{ét}}$}
$(\Gamma=\operatorname{Gal}(\overline K/K))$
can be interpreted as deduced
from the ``geometric'' situation $(\overline X,\overline S)$ over $\overline K$,
or $B_{\overline D^*}\longrightarrow B_{\overline U}$, taking into account
the operations of $\Gamma$ on it. It turns out
that for many questions, it is
this ``geometric'' interpretation (in the
strict sense, i.e. over $\overline K$ algebraically closed) with
operations of a Galois group $\Gamma$, that
is the most convenient.

% L12--L23
If we look at an operation of $\Gamma$
on a topos ($B_{D^*}$ say),
it operates on the discrete topos $B_{\unsure{I|\pi_0}}$
$(I=\pi_0(B_{D^*}))$, and $B_{D^*}\longrightarrow B_I$ is
compatible with the actions of $\Gamma$. \deleted{\unsure{In}}
But the \deleted{structure of} \inserted{topos} $B_{(I,\Gamma)}$ is also
that of $\Gamma$-sets over
$I$ (on which $\Gamma$ operates) --- an induced topos
of $B_\Gamma$ --- \deleted{and the \unsure{fibers over}
an $i\in I$ \unsure{of}} \inserted{its components}
that are connected correspond

% P: handwritten numeral in the lower-left corner; hand not identified.
\margin{174}

\scanpage{176}
% L01--L05
to the orbits of $\Gamma$ on $I$. If $\Gamma$ is
transitive on $I$ \inserted{nonempty} (or if we look at \uline{one}
such orbit...), choosing $i\in I$, the topos
in question identifies with $B_{\Gamma_i}$, where $\Gamma_i$
is the \unsure{Stabilizer} of \deleted{\illegible{1}} $i$ in $\Gamma$.

% L06--L16
\uline{N.B.} Giving an operation of $\Gamma$
on a \uline{discrete} topos $B$ ($=B_I$),
when we interpret it as
a morphism of a
multi-Galois topos $B'=(B,\Gamma)\longrightarrow B_\Gamma$,
is characterized by the fact that
the morphism \inserted{\unsure{of groupoids describing it}}
is \uline{injective} on
arrows, i.e. in terms of a \unsure{system}
\unsure{of homomorphisms} of groups $\Gamma_j\longrightarrow\Gamma$ ($j\in J$
$[\simeq I/\Gamma]$), by the condition that these
homomorphisms be injective. [and then $I$

% L17--L20
is reconstructed as the direct
sum of the $\Gamma/\Gamma_j$ ... ]. Thus a
diagram
\[
 B_{D^*}\longrightarrow B_U
\]

% P: handwritten numeral in the lower-left corner; hand not identified.
\margin{175}

\scanpage{177}
\mspage{88}
% L01--L08
\uline{with action of $\Gamma$} is equivalent to the datum
of a diagram
\[
 B_{D^*,\Gamma}\longrightarrow B_{U,\Gamma}\longrightarrow B_\Gamma
\]
and this is completed (using
the action of $\Gamma$ on $\pi_0(B_{D^*}]$ ---
$B_{D^*}$ being itself deduced from $B_{D^*,\Gamma}\longrightarrow B_\Gamma$
as the inverse image of the
universal torsor) \deleted{\unsure{as|with}} to
\[
 B_{D^*,\Gamma}\longrightarrow B_{I,\Gamma}\longrightarrow B_\Gamma
\]

% L09--L16
 where $B_{D^*,\Gamma}\longrightarrow B_{I,\Gamma}$ can also be
described directly in terms of
$B_{D^*,\Gamma}\longrightarrow B_\Gamma$ (as in an earlier numbe\textsuperscript{r}),
\inserted{at the level of the fundamental groupoids,}
by a passage to the
quotient, \unsure{in|of the} each \unsure{connected component|connected components}
of $\Pi_1 B_{D^*,\Gamma}$, by the kernels of the maps
\unsure{on the homomorphisms deduced|by the homomorphisms deduced} from $\Pi_1 B_{D^*,\Gamma}\longrightarrow\Pi_1 B_\Gamma$.

% L17--L20
--- i.e. this is the canonical factorization
of a homomorphism of \unsure{groupoids|groupoids}
into \unsure{(}\inserted{\unsure{bijective|bijective} (for objects),} epimorphic
\deleted{(for \illegible{1})} \inserted{(for homomorphisms)}\unsure{)} followed by \unsure{monomorphic}
(for homomorphisms). We thus find an

% P: handwritten numeral in the lower-left corner; hand not identified.
\margin{176}

\scanpage{178}
% L01; D01
essentially commutative square
\diagram{\begin{tikzcd}[ampersand replacement=\&]
 \& B_{D^*,\Gamma} \arrow[dl] \arrow[dr] \& \\
 B_{U,\Gamma} \arrow[dr] \& \& B_{I,\Gamma} \arrow[dl] \\
 \& B_\Gamma \&
\end{tikzcd}}

% L02--L03; D02
which \deleted{\unsure{is none other than the one that}}
\inserted{corresponds to the diagram of groupoids}
at the beginning of the notes \deleted{\unsure{that}} \deleted{\unsure{$\Pi$}} \deleted{\unsure{made}} ($\S\,7$)
\diagram{\begin{tikzcd}[ampersand replacement=\&]
 \& \Pi_{D^*} \arrow[dl] \arrow[dr] \& \\
 \Pi_U \arrow[dr] \& \& \Pi_D \arrow[dl] \\
 \& \Pi_e \&
\end{tikzcd}}

% L04--L12
I have the impression of having more or less
understood the mechanism of the actions of
groups on multi-Galois topoi,
and how \deleted{\unsure{this commutes}}
\inserted{the operation of passing}
from a topos $B$ with operation of $\Gamma$
to the ``quotient'' topos $(B,\Gamma)=\text{``}B/\Gamma\text{''}$,
commutes with operations of the type
passing from
\diagram{[$B_{D^*}\longrightarrow B_U$ ; below $B_{D^*}$,
$B_S$, connected to $B_{D^*}$ by a short vertical stroke thicker at its
lower end : \unsure{arrow $B_{D^*}\to B_S$|insertion pointer}.]}
to a $B_{X,S}$, via a
$B_X$ (``plugging the holes''). \deleted{\unsure{I}} The time

% P: handwritten numeral in the lower-left corner; hand not identified.
\margin{177}

\scanpage{179}
\mspage{89}
% L01--L03
seems therefore ripe at last to explain oneself
\unsure{about|with} the hypothetical inverse operation
of ``drilling \unsure{of|of the} holes''.

% L04--L05
\unsure{17.} Supplement on the operations of finite groups
on surfaces. (Supplement to \unsure{$\S\,7$|$\S\,15$})

% L06--L12
Let $U$ be a paracompact connected surface \unsure{!}
such that $\pi_1(U)\ne(1)$ i.e. $U\not\simeq S^2,\mathbb R^2$.
We say that $U$ is ``\unsure{anabelian}'' if
$\pi$ nonabelian (in which case $\operatorname{Centre}(\pi)=1$),
\unsure{i.e.|and} if $U\not\simeq\mathbb C^*,S^1\times S^1$.
Let $\Gamma$
be a finite group operating on $U$. Equivalent
conditions :

% L13--L16
\unsure{b)|a)} \deleted{For} $\forall g\in\Gamma$, $g_U$ is isotopic to the identity.

$b')$ $\forall g\in\Gamma$, $g_U$ is homotopic to the identity.

$c)$ The outer operation of $\Gamma$ on
$\pi_1(U)$ is trivial.

% L17--L20
$d)$ \inserted{(anabelian case)} The extension $\mathcal E$ of $\Gamma$ by $\pi$ is
\unsure{isomorphic} to a product extension, or (abelian case)
\deleted{\illegible{5}} it is central.
\deleted{it \illegible{3} is defined by a}

% L21--L24: heavily superposed revisions at the foot of the page.
\unsure{a)|d)} (anabelian case) $\Gamma$ operates trivially,
(\unsure{abelian|anabelian} case) $\exists$ \deleted{\illegible{3}}
\inserted{structure of a \unsure{topological group}}
\deleted{\unsure{$U\simeq\mathbb C^*,S^1\times S^1$}}
on $U$ (thus $U\simeq\unsure{\mathbb C^*} \unsure{ou} S^1\times S^1$),
\illegible{3}
\deleted{[more intrinsically $U=(\pi\otimes_{\mathbb Z}\unsure{\mathbb C})/\pi$
\illegible{4}]}
\unsure{and a homomorphism of groups $\Gamma\longrightarrow U$, \unsure{operating by} translations}

% M01: upper left margin, written nearly at right angles to the body text.
\margin{That must be the
only counterexample
(in the case \deleted{\illegible{2}}
\unsure{reversal} of an orientation
... ) In any case a
counterexample must be such that
(if $\Gamma$ faithful, $\Gamma\simeq(\pm1)$), \unsure{operates by}
\uuline{anti} \unsure{involution}...}

% M02: oblique note down the left margin; a long return line leads toward M01.
\margin{It is prudent
to assume that
$\Gamma$ operates preserving an orientation
of $U$ (assumed
orientable), otherwise we have
trouble, for example
with \unsure{$z\longrightarrow\overline z^{-1}$} from
$\mathbb C^*$ into $\mathbb C^*$. It
would be necessary to \unsure{clarify|clear up} the
general situation (\unsure{abelian} case)
in the theorem \illegible{4}
case of an anti-involution \unsure{)}.}

% P: handwritten numeral in the lower-left corner; hand not identified.
\margin{178}

\scanpage{180}
% L01
\uline{Proof} $a)\Longrightarrow\unsure{b)}\Longrightarrow b')\Longrightarrow c)$ trivial.

% L02--L08
$c)\Longrightarrow d)$ In anabelian case, this comes
from the fact that $\operatorname{Centre}(\pi)=1$ (and an extension
of $\Gamma$ by $\pi$ is already defined by the outer
action, \deleted{\unsure{taken}} as the inverse image
of the example $1\longrightarrow\pi\longrightarrow\operatorname{Aut}(\pi)
\longrightarrow\operatorname{Autext}\pi\longrightarrow1$ ].
In the abelian case, this is trivial.

% L09--L14
$d)\Longrightarrow a)$ Is the \unsure{not very|not} obvious part. \deleted{\unsure{For the}}
\deleted{In the anabelian case} \unsure{We may assume|Let us assume} that $\Gamma$
operates faithfully, \deleted{to prove that}
\inserted{Anabelian case} \deleted{\illegible{1}}
If we \inserted{\unsure{had|with} $\Gamma\ne1$, let us fix} a splitting of the extension
$\mathcal E$ of $\Gamma$ by $\pi$, we must have
by the theorem of numbe\textsuperscript{r} $15$ $\pi^\Gamma\unsure{=|\simeq}1$, now

% L15--L16
the canonical splitting of $\pi\unsure{\times|\rtimes}\Gamma$ over $\Gamma$
gives $\pi^\Gamma=\pi$, absurd.

% L17--L20
Abelian case $U\simeq\mathbb C^*$ (more intrinsically,
$U$ a torsor under $U_0=\pi\otimes_{\mathbb Z}\mathbb C/\pi$)
or $U\simeq S^1\times S^1$ (more \unsure{intrinsically},
$U$ a torsor under $U_0=\pi\otimes_{\mathbb Z}\mathbb C/\pi$).

% L21--L23
I say that an action of $\Gamma$ on $U$ is
here \unsure{(up to} \deleted{\unsure{isotopy}} \unsure{homeomorphism|conjugation}) defined by an action
of $\unsure{\Gamma|\pi}$ on the torsor --- the class

% M01: note in the left margin, linked by a star and a vertical return line to L21--L23.
\margin{not proved!}

% P: handwritten numeral in the lower-left corner; hand not identified.
\margin{179}